\documentclass[12pt,letterpaper]{article}
\renewcommand{\rmdefault}{cmr}
\renewcommand{\sfdefault}{cmss}
\renewcommand{\ttdefault}{cmtt}
\usepackage[T1]{fontenc}
\usepackage[utf8]{inputenc}
\usepackage[english,spanish]{babel}
\usepackage[letterpaper]{geometry}
\geometry{verbose,nomarginpar,tmargin=3cm,bmargin=3cm,lmargin=3cm,rmargin=3cm}
\setlength{\parskip}{\medskipamount}
\setlength{\parindent}{0pt}
\usepackage{array}
\usepackage{float}
\usepackage{setspace}
\doublespacing

% Paquete de colores para las etiquetas de estado
\usepackage{xcolor}
\definecolor{colorIncluido}{RGB}{29, 78, 216}       % Azul claro / brillante
\definecolor{colorRedactado}{RGB}{30, 58, 138}      % Azul oscuro
\definecolor{colorParcial}{RGB}{107, 114, 128}      % Gris medio
\definecolor{colorConstruccion}{RGB}{156, 163, 175} % Gris claro
\definecolor{colorTextoGris}{RGB}{107, 114, 128}   % Gris notas de figuras/tablas

\usepackage[unicode=true,
 bookmarks=true,bookmarksnumbered=false,bookmarksopen=false,
 breaklinks=false,pdfborder={0 0 0},backref=section,colorlinks=true, linkcolor=blue]
 {hyperref}
\hypersetup{pdftitle={Titulo del proyecto sin tildes},
 pdfauthor={Nombres del los autores sin tildes},
 pdfsubject={tema de la investigacion},
 pdfkeywords={palabras claves}}

\makeatletter

%% Los convertidores HTML no saben que es tabularnewline
\providecommand{\tabularnewline}{\\}

% No generar la salida para \date
\date{}
% Utilizar el punto decimal, en vez de la acostumbrada coma del español.
% Esto está de acuerdo con la Norma NTC/ISO 80000-1
\decimalpoint

% Etiqueta de figura y tabla a la izquierda seguida de punto.
\usepackage[justification=justified, singlelinecheck=false, labelsep = period]{caption}

% Centrar títulos de secciones
\usepackage{sectsty}
\sectionfont{\centering}

% Empezar sección en pagina impar
\let\stdsection\section
\renewcommand\section{\cleardoublepage\stdsection}

% Renombrar las secciones de acuerdo con NTC1486
\addto{\captionsspanish}{\renewcommand*{\contentsname}{\centering CONTENIDO}}
\addto{\captionsenglish}{\renewcommand*{\contentsname}{\centering CONTENIDO}}
\addto{\captionsspanish}{\renewcommand*{\listfigurename}{\centering LISTA DE FIGURAS}}
\addto{\captionsenglish}{\renewcommand*{\listfigurename}{\centering LISTA DE FIGURAS}}
\addto{\captionsspanish}{\renewcommand*{\listtablename}{\centering LISTA DE TABLAS}}
\addto{\captionsenglish}{\renewcommand*{\listtablename}{\centering LISTA DE TABLAS}}
\addto{\captionsspanish}{\renewcommand*{\refname}{BIBLIOGRAF\'IA}}
\addto{\captionsenglish}{\renewcommand*{\refname}{BIBLIOGRAF\'IA}}
\addto{\captionsspanish}{ \renewcommand*{\tablename}{Tabla}}
\addto{\captionsenglish}{ \renewcommand*{\tablename}{Tabla}}
\addto{\captionsenglish}{ \renewcommand*{\figurename}{Figura}}

% Crear un nuevo tipo de sección que sea centrado y sin numeración.
\makeatletter
\newcommand\@csection{\@startsection {section}{1}{\z@}%
  {-3.5ex \@plus -1ex \@minus -.2ex}%
  {2.3ex \@plus.2ex}%
  {\normalfont\Large\bfseries\centering}}
\newcommand{\csection}[1]{%
  \@csection*{#1}%
}
\makeatother

% Centrar la palabra BIBLIOGRAFÍA, sin modificar la entrada en la TOC.
\makeatletter
    \renewenvironment{thebibliography}[1]{%
    \csection{\refname} \bigskip
    \list{\@biblabel{\@arabic\c@enumiv}}
            {\settowidth\labelwidth{\@biblabel{#1}}%
    \leftmargin\labelwidth
    \advance\leftmargin\labelsep
           \@openbib@code
    \usecounter{enumiv}%
    \let\p@enumiv\@empty
    \renewcommand\theenumiv{\@arabic\c@enumiv}}%
    \sloppy
    \clubpenalty4000\@clubpenalty \clubpenalty
    \widowpenalty4000%
    \sfcode`\.\@m}
           {\def\@noitemerr{\@latex@warning{Empty `thebibliography' environment}}%
    \endlist}
\makeatother

% Comando para generar la línea punteada con la etiqueta de estado a la derecha
\newcommand{\tocitem}[3][0pt]{%
  \noindent\hspace{#1}{#2}\ \dotfill\ {#3}\par\vspace{-2pt}%
}

% Marcador visual para secciones pendientes
\newcommand{\enconstruccion}{%
  \vspace{0.5em}%
  \noindent\textcolor{colorConstruccion}{\textit{[En construcción: esta sección se redactará en las próximas sesiones del curso.]}}\par\vspace{1em}%
}


\begin{document}

% ==========================================
% PORTADA
% ==========================================
\begin{titlepage}
    \centering
    \begin{singlespace}
        % 1. ENCABEZADO INSTITUCIONAL
        UNIVERSIDAD PEDAGÓGICA Y TECNOLÓGICA DE COLOMBIA\\
        FACULTAD DE INGENIERÍA\\
        ESCUELA DE INGENIERÍA DE SISTEMAS Y COMPUTACIÓN

        \vfill % Distribución vertical elástica

        % 2. TÍTULO DEL PROYECTO
        \pdfbookmark[1]{Título}{TITULO}
        {\bfseries ANÁLISIS E IDENTIFICACIÓN DE MÉTRICAS AVANZADAS APLICADAS EN LA SELECCIÓN DE ATACANTES DENTRO DE LA PREMIER LEAGUE EN LA TEMPORADA 2025/2026\par}

        \vfill

        % 3. AUTORES

        Joseph Mauricio Fonseca Gonzalez\\


        \vfill

        % 4. CURSO Y CIUDAD/FECHA
        Documento elaborado en el marco del curso Metodología de la Investigación\\[0.5em]
        Tunja, Colombia $\cdot$ 2026
    \end{singlespace}
\end{titlepage}

\singlespacing

% ==========================================
% CONTENIDO (ÍNDICE CON ESTADOS)
% ==========================================
\begin{center}
    {\Large\bfseries CONTENIDO}
\end{center}
\vspace{1em}

\tocitem{\hyperlink{sec:lot}{Índice de tablas}}{}
\tocitem{\hyperlink{sec:lof}{Índice de figuras}}{}
\tocitem{\hyperlink{sec:planteamiento}{\textbf{1. Planteamiento del problema}}}{}
\tocitem[1.5em]{\hyperlink{sec:antecedentes}{1.1 Antecedentes}}{}
\tocitem[1.5em]{\hyperlink{sec:descripcion}{1.2 Descripción de la situación problemática}}{}
\tocitem[1.5em]{\hyperlink{sec:formulacion}{1.3 Formulación del problema}}{}
\tocitem[1.5em]{\hyperlink{sec:consecuencias}{1.4 Consecuencias de no investigar}}{}
\tocitem[1.5em]{\hyperlink{sec:delimitacion}{1.5 Delimitación}}{}
\tocitem{\hyperlink{sec:justificacion}{\textbf{2. Justificación}}}{}
\tocitem{\hyperlink{sec:objetivos}{\textbf{3. Objetivos}}}{}
\tocitem{\hyperlink{sec:alcance}{\textbf{4. Alcance de la investigación}}}{}
\tocitem{\hyperlink{sec:hipotesis}{\textbf{5. Formulación de hipótesis}}}{}
\tocitem{\hyperlink{sec:variables}{\textbf{6. Variables}}}{}
\tocitem{\hyperlink{sec:marco_referencial}{\textbf{7. Marco referencial}}}{}
\tocitem{\hyperlink{sec:diseno_metodologico}{\textbf{8. Diseño metodológico}}}{}
\tocitem{\hyperlink{sec:resultados}{\textbf{9. Resultados y discusión}}}{}
\tocitem{\hyperlink{sec:conclusiones}{\textbf{10. Conclusiones y recomendaciones}}}{}
\tocitem{\hyperlink{sec:referencias}{Referencias}}{}
\tocitem{\hyperlink{sec:anexos}{Anexos}}{}

\vspace{2.5em}
% ==========================================
% ÍNDICE DE TABLAS
% ==========================================
\begin{center}
    {\Large\bfseries ÍNDICE DE TABLAS}
\end{center}
\vspace{0.5em}
\begin{quote}
    \centering\color{colorTextoGris}\textit{Aún no hay tablas incorporadas al cuerpo del documento. Se listan aquí a medida que cada fase incorpore tablas propias (matrices, resultados, operacionalización de variables, etc.).}
\end{quote}

\vspace{1.5em}

% ==========================================
% ÍNDICE DE FIGURAS
% ==========================================
\begin{center}
    {\Large\bfseries ÍNDICE DE FIGURAS}
\end{center}
\vspace{0.5em}
\begin{quote}
    \centering\color{colorTextoGris}\textit{Aún no hay figuras incorporadas al cuerpo del documento. Se listan aquí a medida que cada fase incorpore figuras propias (diagramas, gráficos, esquemas, etc.).}
\end{quote}

\vspace{2.5em}
\onehalfspacing

% ==========================================
% CUERPO DE LA INVESTIGACIÓN
% ==========================================
\hypertarget{sec:planteamiento}{}%
\begin{center}
    {\Large\bfseries 1. PLANTEAMIENTO DEL PROBLEMA}
\end{center}
\vspace{1em}

\hypertarget{sec:antecedentes}{}%
\noindent{\large\bfseries 1.1 ANTECEDENTES}\par\vspace{0.5em}

El uso del análisis cuantitativo aplicado a la toma de decisiones en el deporte profesional se originó con la sabermetría implementada por los Oakland Athletics en la Major League Baseball (MLB) durante la temporada 2002 [1]. Con un presupuesto de 44 millones de dólares, equivalente a un tercio del capital de franquicias como los New York Yankees (125 millones de dólares), la institución sustituyó los criterios de observación tradicional por modelos estadísticos para clasificar y seleccionar jugadores en todas las posiciones, logrando clasificar a los playoffs en 2002 y 2003 [1]. 

Este antecedente demostró que las ineficiencias del mercado deportivo podían ser explotadas mediante la cuantificación del rendimiento, sentando las bases para la posterior adopción de datos en la gestión integral de plantillas de fútbol [2].

En el fútbol europeo, esta aproximación tuvo su primera aplicación formal en el Liverpool FC entre 2010 y 2012 bajo la dirección de Damien Comolli [3]. El enfoque priorizó la detección de ineficiencias tácticas y funcionales en las distintas líneas del terreno de juego (defensa, creación y finalización) para incorporar futbolistas que solventaran las carencias estructurales específicas del plantel [3]. Aunque esta experiencia inicial en un club de élite tuvo corta duración, estableció la viabilidad de utilizar indicadores cuantitativos para la planificación global de plantillas[4].

A partir de este antecedente, la visión analítica fue perfeccionada por instituciones que buscaban competir contra clubes de mayor capacidad financiera, destacando los casos del Brighton \& Hove Albion y el Brentford FC [5]. Tras asumir el control de sus respectivos clubes en 2009 y 2012, Tony Bloom y Matthew Benham estructuraron modelos algorítmicos aplicados orientados a identificar futbolistas con una relación entre rendimiento-costo efectiva [5]. La aplicación de estas estructuras en la confección del equipo impulsó el ascenso de ambos clubes desde la tercera categoría inglesa hasta la Premier League (en 2016 y 2021), consolidando proyectos competitivos y sostenibles en la élite [6].

El éxito de estos modelos aceleró la incorporación de departamentos de datos en las cinco principales ligas europeas (Premier League, LaLiga, Serie A, Bundesliga y Ligue 1), donde el uso de estructuras analíticas se ha masificado en la última década [7]. 

Esta expansión ha sido impulsada por la acelerada inflación en los precios de traspaso y por el incremento en los ingresos televisivos, permitiendo a los clubes adquirir tecnologías de seguimiento (tracking data) y de eventos (event data) para medir cuantitativamente cada fase del juego y cada rol dentro del campo [7].

En paralelo, la literatura académica sobre analítica deportiva en el fútbol ha crecido de forma desproporcionada hacia la modelación matemática de eventos aislados o el estudio de clubes de élite con infraestructuras propietarias de alto costo [8]. 

Existe, sin embargo, una escasa literatura científica que documente de manera transparente cuáles son las métricas específicas y las ponderaciones reales que utilizan los clubes ingleses al momento de evaluar e incorporar futbolistas en sus plantillas [8]. 

Los estudios disponibles suelen analizar indicadores teóricos en entornos ideales, advirtiendo que sus hallazgos no son fácilmente extrapolables a los procesos reales de selección y scouting que ocurren en el mercado de fichajes [8].

Una revisión exploratoria inicial realizada por el equipo sobre una muestra de 15 informes y publicaciones de scouting enfocado en plantillas de la Premier League (temporadas 2023-2025) evidenció que en 12 de ellos se reportaban métricas heterogéneas sin una definición metodológica estandarizada sobre qué indicadores determinan la selección de un jugador para una posición específica [9]. 

Aunque esta revisión posee un carácter exploratorio y no constituye una muestra representativa de todo el mercado, actúa como una señal inicial que confirma la opacidad y la escasez de documentación en la literatura sobre los criterios métricos reales empleados en el fútbol inglés [9].

Esta falta de literatura consolidada sobre las métricas de selección expone a las instituciones de menores recursos y a los analistas independientes a un escenario de alta incertidumbre, favoreciendo decisiones basadas en muestras pequeñas de partidos o en apreciaciones subjetivas [10].

Estos elementos (la evolución del análisis cuantitativo para la conformación de plantillas, la masificación de datos en las grandes ligas y la escasa literatura documentada sobre las métricas reales empleadas por los clubes ingleses) describen el contexto del problema, pero no permiten, por sí solos, establecer qué métricas avanzadas específicas y qué ponderaciones contextuales utilizan o deberían utilizar los clubes de la Premier League para la evaluación y fichaje de sus futbolistas en la temporada 2025/2026 (y, de manera particular, en roles de alto impacto como los atacantes): para eso hace falta un análisis directo sobre los datos de dicha temporada que permita caracterizar sistemáticamente estos criterios de selección. Los datos históricos y reportes citados documentan la situación que motivó esta propuesta; no sustituyen el estudio empírico planeado, el cual se ejecutará sobre el registro de datos de la temporada 2025/2026.

\vspace{1em}

\hypertarget{sec:descripcion}{}%
\noindent{\large\bfseries 1.2 DESCRIPCIÓN DE LA SITUACIÓN PROBLEMÁTICA}\par
Como se documentó en los antecedentes (numeral 1.1), el uso de la analítica de datos en el fútbol europeo ha crecido aceleradamente, pero la toma de decisiones sobre la selección y fichaje de atacantes sigue dominada por métricas tradicionales reduccionistas o por la descontextualización de indicadores avanzados ($xG$, $xA$). Actualmente no existe, en la literatura accesible ni en la práctica de clubes de menor presupuesto, evidencia o un marco metodológico estandarizado que permita determinar con rigor qué métricas abiertas y qué ponderaciones tácticas contextuales diferencian objetivamente a un atacante en la Premier League durante la temporada 2025/2026: la revisión exploratoria inicial reveló que, en 12 de 15 informes analizados, las evaluaciones ignoraban el modelo de juego del equipo de origen y se limitaban al conteo bruto de goles o a muestras pequeñas de partidos [9]. Esta ausencia de herramientas accesibles afecta directamente la planificación deportiva de las instituciones con recursos limitados, obligando a que las decisiones de inversión en contrataciones y la promoción de talentos se tomen por intuición, bajo la influencia de agentes o basadas en apreciaciones subjetivas [10].

\vspace{1em}

\hypertarget{sec:formulacion}{}%
\noindent{\large\bfseries 1.3 FORMULACIÓN DEL PROBLEMA}\par
A partir del planteamiento anterior, la pregunta de investigación que orienta este estudio es: ¿de qué manera el análisis e integración de métricas avanzadas de rendimiento accesible permite caracterizar y diferenciar perfiles de atacantes de forma consistente en la Premier League durante la temporada 2025/2026, en comparación con las evaluaciones basadas en métricas tradicionales?

Esta pregunta se delimita a la temporada 2025/2026 de la Premier League, a futbolistas de perfil ofensivo (atacantes: delanteros centros, extremos e interiores de ataque), y al uso de datos cuantitativos de eventos de acceso abierto o bajo costo. Se sistematiza en las siguientes sub-preguntas: (a) ¿cuáles métricas avanzadas accesibles presentan mayor capacidad de diferenciación en la efectividad y contribución táctica de los atacantes evaluados?; (b) ¿cómo altera el contexto colectivo del equipo (como el porcentaje de posesión y volumen de generación) la ponderación de dichos indicadores avanzados?; y (c) ¿existen diferencias en la caracterización funcional entre atacantes con registros tradicionales (goles y asistencias) similares que deban ajustarse mediante este modelo?

\vspace{1em}

\hypertarget{sec:consecuencias}{}%
\noindent{\large\bfseries 1.4 CONSECUENCIAS DE NO INVESTIGAR}\par
Si esta situación persiste, es probable que las instituciones de recursos limitados y los analistas independientes sigan evaluando y contratando futbolistas de perfil ofensivo sin criterio metodológico explícito, con dos riesgos simultáneos que impactan la sostenibilidad del club: asumir contrataciones a sobreprecio de jugadores cuyo rendimiento decae al cambiar de sistema táctico (comprometiendo el capital en salarios de alto costo) o descartar talentos accesibles y efectivos por no figurar en conteos brutos tradicionales. Cada ventana de transferencias que transcurre sin un marco metodológico abierto y contextualizado es, previsiblemente, un periodo más de política de fichajes basada en intuición, especulación comercial y sesgos de evaluación.

\vspace{1em}

\hypertarget{sec:delimitacion}{}%
\noindent{\large\bfseries 1.5 DELIMITACIÓN}\par
Este estudio se delimita a la caracterización y diferenciación de perfiles de atacantes mediante métricas avanzadas accesibles y su contexto colectivo —no al desarrollo de una plataforma de software comercial ni a modelos de aprendizaje automático predictivos de lesiones—, en futbolistas con minutos jugados en la Premier League durante la temporada 2025/2026, utilizando datos cuantitativos de eventos de acceso público o bajo costo. La ejecución y análisis se realizan sobre el registro de datos de la temporada 2025/2026, no sobre periodos o ligas anteriores.

Queda fuera del alcance el seguimiento biométrico y de posicionamiento en tiempo real (optical tracking data de alto costo) —filtrado por no ser viable de adquirir o procesar sin infraestructura de élite—, así como el análisis de posiciones defensivas o de portería, la predicción de valor económico de mercado en dinero, el impacto de factores psicológicos individuales y la evaluación de competencias fuera del contexto inglés.
\begin{center}
    {\Large\bfseries 2. JUSTIFICACIÓN}
\end{center}
\enconstruccion

\begin{center}
    {\Large\bfseries 3. OBJETIVOS}
\end{center}
\enconstruccion
\begin{center}
    {\Large\bfseries 4. ALCANCE DE LA INVESTIGACIÓN}
\end{center}

Este estudio es de tipo explicativo: busca establecer si la integración e interpretación contextualizada de métricas avanzadas de rendimiento accesible (variable independiente / causa) produce una diferencia medible en la precisión de la caracterización funcional y en la reducción de la incertidumbre al diferenciar perfiles de atacantes (variables dependientes / efectos), en comparación con un modelo de evaluación tradicional basado en estadísticas brutas (goles y asistencias). La elección de este alcance se justifica en que la investigación no se limita a describir el uso de datos en el fútbol ni a observar la simple correlación entre variables, sino que pone a prueba el impacto de un marco de ponderación táctico-contextual sobre la evaluación de jugadores. Para ello, se apoya en el diagnóstico descriptivo previo sobre la toma de decisiones en el *scouting* (numeral 1) y en los estudios correlacionales existentes sobre la efectividad de métricas como goles esperados ($xG$) y asistencias esperadas ($xA$) en clubes de élite, evaluando si esta metodología permite diferenciar objetivamente a futbolistas con registros tradicionales similares en la Premier League durante la temporada 2025/2026.

Limitaciones de la investigación

Al basarse en datos secundarios de eventos (*event data*) de acceso abierto y no en un entorno controlado con datos de seguimiento óptico en tiempo real (*optical tracking data*), las inferencias sobre comportamientos tácticos sin balón deben asumirse con mayor cautela que las derivadas de mediciones de alta resolución de clubes de élite; la muestra se limita de manera exclusiva a futbolistas de perfil ofensivo con minutos jugados en la Premier League durante la temporada 2025/2026, lo que restringe qué tanto se pueden generalizar o transferir directamente las ponderaciones cuantitativas a otras ligas continentales, divisiones inferiores o demarcaciones defensivas sin una recalibración previa de los parámetros competitivos; y el modelo de ponderación multicriterio diseñado para matizar las métricas según el contexto del equipo se valida dentro de este mismo estudio, por lo que sus propiedades metodológicas (consistencia interna, sensibilidad del modelo y nivel de concordancia) deben reportarse empíricamente junto con los resultados finales del proyecto, no darse por supuestas previamente.


\hypertarget{sec:hipotesis}{}%
\begin{center}
    {\Large\bfseries 5. FORMULACIÓN DE HIPÓTESIS}
\end{center}
\vspace{1em}

\noindent\textbf{H1a (hipótesis de trabajo --- caracterización funcional):} Los atacantes evaluados mediante el modelo de métricas avanzadas ajustadas por contexto táctico presentan una mayor precisión en la diferenciación de su perfil funcional que cuando son evaluados mediante el modelo tradicional basado en estadísticas basicas (goles y asistencias), en futbolistas con registros estadísticos equivalentes.\par\vspace{0.8em}

\noindent\textbf{H0a (hipótesis nula --- caracterización funcional):} No existe diferencia en la precisión de la diferenciación del perfil funcional de los atacantes entre la evaluación realizada con el modelo de métricas avanzadas ajustadas por contexto táctico y la realizada con el modelo tradicional basado en estadísticas basicas.\par\vspace{1.2em}

\noindent\textbf{H1b (hipótesis de trabajo --- incertidumbre en la evaluación):} La aplicación del modelo de métricas avanzadas contextualizadas presenta un menor grado de incertidumbre y discrepancia en la clasificación del rendimiento real de los atacantes que la evaluación basada exclusivamente en estadísticas tradicionales basicas.\par\vspace{0.8em}

\noindent\textbf{H0b (hipótesis nula --- incertidumbre en la evaluación):} No existe diferencia en el grado de incertidumbre ni en la clasificación del rendimiento real de los atacantes entre la aplicación del modelo de métricas avanzadas contextualizadas y la evaluación basada exclusivamente en estadísticas tradicionales basicas.
\begin{center}
    {\Large\bfseries 6. VARIABLES}
\end{center}
\enconstruccion

\begin{center}
    {\Large\bfseries 7. MARCO REFERENCIAL}
\end{center}
\enconstruccion

\begin{center}
    {\Large\bfseries 8. DISEÑO METODOLÓGICO}
\end{center}
\enconstruccion
\begin{center}
    {\Large\bfseries 9. RESULTADOS Y DISCUSIÓN}
\end{center}
\enconstruccion

\begin{center}
    {\Large\bfseries 10. CONCLUSIONES Y RECOMENDACIONES}
\end{center}
\enconstruccion

% ==========================================
% REFERENCIAS
% ==========================================
\renewcommand{\refname}{REFERENCIAS}
\begin{thebibliography}{9}
\bibitem{ref1}
A. Smith, ``Análisis del rendimiento académico en cursos introductorios de programación,'' \textit{Revista Iberoamericana de Educación en Ingeniería}, vol. 12, no. 2, pp. 45--58, 2023.
\end{thebibliography}

% ==========================================
% ANEXOS
% ==========================================
\begin{center}
    {\Large\bfseries ANEXOS}
\end{center}
\enconstruccion

\end{document}